% !TeX root = ../main.tex


\chapter{Introduction}
\label{ch:intro}


\section{Context and motivation}
\label{sec:intro:context}

Location-based social networks record visits through \emph{check-ins}. Each check-in
states that a user visited a place, or point of interest (POI), at a given time.
Although these traces are noisy, human movement is highly regular: an analysis of
large-scale mobility traces estimated the potential predictability of an individual's
next location at about 93 percent~\cite{song2010limits}. Anticipating what kind of
place a person will visit and where that visit will occur can support recommendation,
navigation, transit planning, and the allocation of resources by area. The broader
study of mobility also informs urban planning, disease-spread analysis, and pollution
research~\cite{luca2021mobilitysurvey}, while place categories provide the semantic
information used by location-based services~\cite{Xu2023}. Some methods used in this
dissertation also come from neighboring geospatial tasks. In particular, the spatial
encoders examined in the second study were first validated for applications such as
species recognition and remote-sensing
classification~\cite{mai2023sphere2vecgeneralpurposelocationrepresentation,wu2024torchspatial}.

This dissertation studies two properties of the next visit. \emph{Next-category
prediction} identifies the semantic category of that visit, and \emph{next-region
prediction} identifies where it will occur. These tasks differ from \emph{next-place
prediction}, which identifies the exact establishment and is outside the scope of
this work. The first two studies also use \emph{category classification}, a static
task that predicts the category of a POI held out from classifier training. The task
pair changes in the final study, as explained in Section~\ref{sec:intro:arc}.
Chapter~\ref{ch:fundamentals} formally defines the three tasks that appear in the
experiments.

Because next category and next region are predicted from the same visit history, one
model may serve both tasks. Multitask learning (MTL) trains related tasks jointly so
that they can share information~\cite{caruana1997multitask}. This approach has been
used in vision~\cite{kokkinos2016ubernet}, clinical
prediction~\cite{lipton2015learning}, and language
modeling~\cite{wei2022finetuned}. For the setting studied here, its operational appeal
is a single model to maintain and one forward pass that produces both predictions,
instead of two dedicated single-task models. Joint training does not guarantee better
predictions, however. Shared parameters can harm one task, a failure known as negative
transfer. Whether MTL would help these POI prediction tasks, and which design choices
would determine the result, was an open question at the start of this research.


\section{Research question}
\label{sec:intro:arc}

\textbf{Does multitask learning help point-of-interest prediction (next category and
next region), and what does the answer depend on?}

The three studies answer this question in sequence. The first reports a negative
result, the second identifies its main bottleneck, and the third tests the resulting
solution. This progression is itself part of the contribution because each study
narrows the explanation supported by the evidence.

Chapter~\ref{ch:cbic}, published at CBIC 2025, introduces MTLnet, the first joint model
developed in this research. It combines a place-level graph embedding with hard
parameter sharing and predicts static place categories and the next category in a
sequence. For that configuration, the joint model did not consistently outperform
the two dedicated single-task models and required more training time. The study
treated this null result as a finding and proposed three possible explanations,
including an input representation that might not provide enough information for both
tasks.

Chapter~\ref{ch:courb}, published at CoUrb 2026, examines that explanation through a
controlled comparison. The study keeps MTLnet unchanged and replaces its monolithic
64-dimensional place embedding with separate spatial, temporal, and categorical
encoders. Category performance then increases substantially in every state evaluated.
Because the architecture remains fixed, the study identifies the input representation
as the bottleneck for that stage of the research.

Chapter~\ref{ch:mobiwac}, submitted to MobiWac 2026 and under review, develops the
resolution. Any place-level embedding assigns a place the same vector on every visit; a
representation that cannot tell a weekday lunch from a Saturday night out at the same
place is working against both tasks at once. That observation is the hypothesis the
final study tests. It moves the representation to the check-in level: one vector per
visit. Acting on another of the first study's candidate explanations, it also replaces
the shared hidden layers with cross-attention between task-specific streams, addressing
the earlier concern about hard parameter sharing. Under a check-in level representation, static
category classification is a less natural companion task than a second sequential target,
so the final task pair becomes next category and next region. Across five United States
states and Istanbul, the study compares the joint and dedicated models with paired superiority
and non-inferiority tests. The results support a positive answer for this configuration;
Chapter~\ref{ch:mobiwac} reports the task-specific outcomes.

The task pair thus changes across the studies: the first two combine static category
classification with next-category prediction, whereas the final study combines
next-category and next-region prediction. The dissertation does not treat these as one
unchanged experiment. It presents a sequence in which a result valid for one
configuration leads to a diagnosis and then to a different, explicitly bounded
solution.


\section{Objectives}
\label{sec:intro:objectives}

The general objective is to determine whether multitask learning helps next-category
and next-region prediction and to identify the conditions that shape the answer. Four
specific objectives connect this question to the article chapters and to their
consolidation:

\begin{enumerate}
    \item Evaluate whether a joint model with hard parameter sharing benefits static
    category classification and next-category prediction when compared with
    dedicated single-task models (Chapter~\ref{ch:cbic}).
    \item Determine how the input representation affects that result by enriching and
    decomposing the input while keeping the architecture fixed
    (Chapter~\ref{ch:courb}).
    \item Design and validate a check-in-level representation and a joint model for
    next-category and next-region prediction (Chapter~\ref{ch:mobiwac}).
    \item Consolidate the evidence under the user-disjoint cross-validation,
    significance-testing, and non-inferiority protocol used in the final study
    (Chapter~\ref{ch:conclusion}).
\end{enumerate}


\section{Scope and assumptions}
\label{sec:intro:scope}

The experiments use observational check-in traces from two sources: Gowalla data for
five United States states (Alabama, Arizona, Florida, California, Texas) and the Istanbul subset of the Massive-STEPS benchmark. The prediction targets are the next category, over the
seven-class taxonomy defined in Chapter~\ref{ch:fundamentals}, and the next region, over census tracts in
the United States and mahalles in Istanbul.
The joint model operates under a single-model constraint: one trained artifact must produce
both outputs in one forward pass. These choices define the design scope.
Chapter~\ref{ch:fundamentals} specifies the output spaces and geographic units, and
Chapter~\ref{ch:conclusion} discusses the limitations that affect how the results may
generalize.


\section{Organization of this dissertation}
\label{sec:intro:organization}

This dissertation is organized as a collection of works, where each central chapter is an
independent article. The articles were re-typeset in a common format, and the second
article was translated into English. Their content follows the published versions or,
for the article under review, the submitted manuscript.
% Corrections introduced during adaptation are recorded in the errata appendix of \extravolume.

\begin{itemize}
    \item \textbf{Chapter~\ref{ch:fundamentals}} defines the shared background: POI
    prediction tasks, mobility representations, multitask learning, datasets,
    metrics, and evaluation protocols.
    \item \textbf{Chapter~\ref{ch:cbic}} reproduces the first article, published in
    English at the XVII Congresso Brasileiro de Intelig\^encia Computacional
    (CBIC 2025, DOI~10.21528/CBIC2025-1191324). This dissertation's author is its
    first author.
    \item \textbf{Chapter~\ref{ch:courb}} presents an English translation of the second
    article, published in Portuguese at the X Workshop de Computa\c{c}\~ao Urbana
    (CoUrb 2026, DOI~10.5753/courb.2026.22960), held with the Simp\'osio
    Brasileiro de Redes de Computadores e Sistemas Distribu\'idos (SBRC 2026).
    Tarik S.\ Paiva is the first author. This dissertation's author is the second
    author, contributed the MTLnet baseline, and presented the work at the event.
    \item \textbf{Chapter~\ref{ch:mobiwac}} presents the third article, submitted to the
    23rd ACM International Symposium on Mobility Management and Wireless Access
    (MobiWac 2026) and under review. This dissertation's author is its first
    author.
    \item \textbf{Chapter~\ref{ch:conclusion}} answers the research question across the
    three studies, states their joint limitations, and connects each limitation
    to future work.
\end{itemize}

Each article chapter begins with a preface that identifies its venue and publication
status. The prefaces also distinguish the conclusion supported at the time of each
article from the final position of the dissertation.


\section{Contributions}
\label{sec:intro:contributions}

The contributions are organized into four groups.

\begin{description}
    \item[Theoretical.] The studies show that the input representation and its sharing
    topology determine whether MTL helps these POI prediction tasks. A null result
    with a place-level representation and hard sharing therefore does not conflict
    with a positive result obtained from a check-in-level representation and a
    different form of sharing.
    \item[Software.] The work provides MTLnet (Chapter~\ref{ch:cbic}), Check2HGI and the
    joint model built on it (Chapter~\ref{ch:mobiwac}), and reproducible training
    and evaluation pipelines for Chapters~\ref{ch:cbic} and~\ref{ch:mobiwac}. The
    corresponding code repositories are referenced in those chapters.
    \item[Empirical.] The final study compares joint and dedicated models on six
    datasets under user-disjoint cross-validation. Each configuration has twenty
    fitted models: four repetitions of the complete five-fold experiment, each
    with a different random initialization and the same fixed folds. The analysis
    uses paired tests on the four per-seed means, non-inferiority tests, and a
    leakage audit. Because every seed uses the same fold partition, the reported
    intervals do not include uncertainty from resampling the user splits.
    \item[Practical.] One model produces both predictions in one forward pass. It
    outperforms the dedicated model on next category and either outperforms or
    matches it on next region under the statistical criteria defined in this
    dissertation.
\end{description}
